\section{Glosario}
\begin{itemize}
    \item \textbf{Adapter (Adaptador).} Patrón de diseño GoF que traduce la interfaz de un componente hacia otra que el sistema espera. Cada pasarela de pago tiene su propio Adapter (WompiAdapter, RapydAdapter, etc.) que traduce el contrato genérico del SDK hacia las convenciones específicas de esa pasarela.
    \item \textbf{API REST.} Interfaz de programación que expone operaciones a través de solicitudes HTTP (GET, POST, etc.) siguiendo el estilo arquitectónico REST.
    \item \textbf{API de Simulación.} Servicio web desarrollado en este proyecto que replica el comportamiento de las cuatro pasarelas de pago soportadas, permitiendo probar integraciones sin usar credenciales reales.
    \item \textbf{ASR (Architecturally Significant Requirement).} Requisito Arquitectónicamente Significativo: requisito funcional o no funcional que tiene un impacto directo y relevante sobre las decisiones de diseño de la arquitectura del sistema.
    \item \textbf{Backoff exponencial.} Estrategia de reintento en la que el tiempo de espera entre intentos sucesivos aumenta de forma exponencial (por ejemplo, 1s, 2s, 4s), utilizada para evitar sobrecargar un servicio que está fallando temporalmente.
    \item \textbf{CBO (Coupling Between Objects).} Métrica del conjunto de Chidamber y Kemerer (1994) que mide el número de clases a las que una clase está acoplada. Se usa en este proyecto para comparar el acoplamiento resultante de integrar pasarelas con y sin el framework.
    \item \textbf{Credenciales de sandbox.} Claves o tokens de prueba proporcionados por una pasarela de pago para realizar transacciones simuladas sin mover dinero real.
    \item \textbf{CK (Métricas de Chidamber y Kemerer).} Conjunto de métricas de calidad de código orientado a objetos (entre ellas WMC, CBO y RFC), propuesto por Chidamber y Kemerer en 1994, usado en este proyecto para evaluar objetivamente el impacto del framework.
    \item \textbf{DDD (Domain-Driven Design).} Enfoque de diseño de software que organiza el sistema alrededor de un modelo del dominio del negocio, distinguiendo entre Entidades, Objetos de Valor y Servicios de Dominio, tal como se aplica en el modelo de dominio de este documento (sección 3).
    \item \textbf{Entidad (DDD).} En Domain-Driven Design, un objeto que tiene una identidad propia que persiste a lo largo del tiempo, independientemente de sus atributos. Transaction es la única entidad del dominio en este proyecto.
    \item \textbf{Facade (Fachada).} Patrón de diseño GoF que provee una interfaz simplificada y unificada hacia un subsistema más complejo. PaymentFacade es el único punto de entrada que el desarrollador consume del SDK.
    \item \textbf{Factory (Fábrica).} Patrón de diseño GoF que encapsula la lógica de creación de objetos, ocultando el proceso de instanciación del código cliente. GatewayFactory retorna el Adapter correspondiente a la pasarela activa.
    \item \textbf{Fastify.} Framework web para Node.js orientado al alto rendimiento, utilizado para implementar la API de Simulación.
    \item \textbf{FURPS+.} Modelo de clasificación de atributos de calidad de software (Functionality, Usability, Reliability, Performance, Supportability, más restricciones adicionales), usado en este documento para organizar los atributos de calidad de la arquitectura.
    \item \textbf{Gateway (Pasarela de pago).} Servicio externo que procesa transacciones de pago en nombre de un comercio. Las cuatro pasarelas soportadas en este proyecto son Wompi, Rapyd, Mercado Pago y Kushki.
    \item \textbf{HMAC-SHA256.} Algoritmo de código de autenticación de mensajes basado en la función hash SHA-256, usado por Mercado Pago y Kushki para firmar sus webhooks y permitir verificar que no fueron alterados.
    \item \textbf{Hexagonal, Arquitectura (Puertos y Adaptadores).} Estilo arquitectónico que separa el núcleo de negocio de un sistema de sus dependencias externas mediante puertos (interfaces) y adaptadores (implementaciones concretas), adoptado como base del SDK en el ADR-01.
    \item \textbf{HTTPS.} Protocolo HTTP sobre una capa de cifrado TLS, usado obligatoriamente en toda comunicación del framework con las pasarelas de pago y la API de Simulación.
    \item \textbf{Idempotencia.} Propiedad de una operación que produce el mismo resultado sin efectos adicionales, sin importar cuántas veces se ejecute. Es relevante para el manejo de reintentos y de webhooks duplicados, un escenario que cualquier pasarela puede generar y que el Webhook Verifier debe manejar sin asumir que cada notificación llega exactamente una vez.
    \item \textbf{JSON Schema.} Especificación para describir y validar la estructura de documentos JSON, utilizada por Fastify en la API de Simulación para validar solicitudes y respuestas de cada escenario.
    \item \textbf{Mock (payload simulado).} Objeto o respuesta construida artificialmente para replicar, con fines de prueba, el comportamiento de un componente real; en este caso, las respuestas nativas de cada pasarela de pago.
    \item \textbf{npm.} Gestor de paquetes de Node.js utilizado para distribuir e instalar el SDK del framework mediante un único comando.
    \item \textbf{Objeto de Valor (Value Object).} En Domain-Driven Design, un objeto sin identidad propia que se define únicamente por sus atributos y que debería ser inmutable. Amount, Currency y Payer, entre otros, son objetos de valor en este proyecto.
    \item \textbf{OpenAPI 3.0.} Especificación estándar para describir APIs REST de forma legible tanto por humanos como por máquinas, expuesta públicamente por la API de Simulación en el endpoint /docs.
    \item \textbf{PaaS (Platform as a Service).} Modelo de servicio en la nube que provee una plataforma completa (runtime, infraestructura) para desplegar aplicaciones sin gestionar servidores directamente. La API de Simulación se despliega sobre una plataforma PaaS (Render).
    \item \textbf{PCI DSS (Payment Card Industry Data Security Standard).} Estándar de seguridad para el manejo de datos de tarjetas de pago. El framework delega su cumplimiento a las pasarelas integradas, ya que no procesa ni almacena datos de tarjeta directamente.
    \item \textbf{Puerto (Port).} En la Arquitectura Hexagonal, una interfaz definida por el núcleo del sistema que especifica qué operaciones necesita del exterior, sin conocer cómo se implementan. PaymentGatewayPort es el puerto que todos los Adapters de pasarela deben implementar.
    \item \textbf{RFC (Response For a Class).} Métrica del conjunto de Chidamber y Kemerer que mide el número de métodos que pueden ejecutarse en respuesta a un mensaje recibido por una clase, incluyendo los suyos propios y los que invoca.
    \item \textbf{RF (Requisito Funcional).} Requisito que describe una funcionalidad concreta que el sistema debe realizar, numerado en este proyecto como RF-01, RF-02, etc.
    \item \textbf{Sandbox.} Entorno de pruebas aislado que provee una pasarela de pago para simular transacciones sin procesar dinero real.
    \item \textbf{SDK (Software Development Kit).} Conjunto de herramientas, en este caso una biblioteca de código distribuida como paquete npm, que un desarrollador integra en su aplicación para consumir las capacidades de pago del framework.
    \item \textbf{Servicio de Dominio.} En Domain-Driven Design, una operación de negocio que no pertenece naturalmente a ninguna Entidad ni Objeto de Valor por no tener datos propios que persistir. WebhookVerifier es un servicio de dominio en este proyecto.
    \item \textbf{SHA-256.} Función hash criptográfica utilizada, entre otros usos, por Wompi para firmar sus webhooks.
    \item \textbf{TypeScript.} Superconjunto tipado de JavaScript, lenguaje en el que está implementado el SDK y la API de Simulación de este proyecto.
    \item \textbf{Webhook.} Mecanismo mediante el cual un sistema notifica a otro, de forma asíncrona, sobre la ocurrencia de un evento, enviando una solicitud HTTP POST hacia una URL previamente configurada. Las pasarelas de pago lo usan para notificar cambios de estado en una transacción.
    \item \textbf{WMC (Weighted Methods per Class).} Métrica del conjunto de Chidamber y Kemerer que mide la complejidad de una clase como la suma de la complejidad de sus métodos.
\end{itemize}

